%! AP Statistics — Unit 4, Lesson 4.3 — Guided Notes KEY
\documentclass[10pt]{article}
\usepackage{apstats-article}
\usepackage{apstats-key}

\begin{document}

\pageheader{Unit 4, Lesson 4.3}{Guided Notes --- KEY}
\namedateperiod

\vspace{0.12in}

\begin{objectivebox}
\small By the end of this lesson, I will be able to\dots\\[4pt]
\ansline{Calculate a probability using a sample space and the classical definition.}\\[1pt]
\ansline{Apply the complement rule: $P(E^c) = 1 - P(E)$.}\\[1pt]
\ansline{Interpret a probability as a long-run relative frequency.}
\end{objectivebox}

\vspace{0.1in}

\begin{vocabbox}
\noindent\textbf{\textcolor{navy}{Sample space ($S$):}}\\[1pt]%
\ansline{The set of all possible non-overlapping outcomes of a random process.}\\[1pt]\ansline{}\\[3pt]
\noindent\textbf{\textcolor{navy}{Event:}}\\[1pt]%
\ansline{An outcome or set of outcomes (a subset of the sample space).}\\[1pt]\ansline{}\\[3pt]
\noindent\textbf{\textcolor{navy}{Equally likely outcomes:}}\\[1pt]%
\ansline{Every outcome in the sample space has the same probability.}\\[1pt]\ansline{}\\[3pt]
\noindent\textbf{\textcolor{navy}{Classical probability:}}\\[1pt]%
\ansline{$P(E) = $ (number of outcomes in $E$) / (total number of outcomes in $S$),}\\[1pt]\ansline{used when all outcomes are equally likely.}\\[3pt]
\noindent\textbf{\textcolor{navy}{Complement ($E^c$):}}\\[1pt]%
\ansline{The event that $E$ does NOT occur: all outcomes in $S$ not in $E$.}\\[1pt]\ansline{}\\[3pt]
\noindent\textbf{\textcolor{navy}{Complement rule:}}\\[1pt]%
\ansline{$P(E^c) = 1 - P(E)$}\\[1pt]\ansline{}\\[3pt]
\end{vocabbox}

\vspace{0.1in}

\begin{hookbox}
\small A bag holds 3 red and 5 blue poker chips.
\begin{itemize}[itemsep=3pt]
  \item List the sample space for one draw: $S = \{$\ \ans{R1, R2, R3, B1, B2, B3, B4, B5}\ $\}$
  \item How many total outcomes? \ans{8}
  \item How many outcomes are ``red''? \ans{3}
  \item Probability of drawing red: $P(\text{red}) = $ \ans{$3/8$} $=$ \ans{$0.375$}
\end{itemize}
\end{hookbox}

\vspace{0.1in}

\begin{notesbox}{The Classical Definition of Probability}
\small
\textbf{VAR-4.A.1 --- Sample Space:} The sample space is the set of all possible
\ans{non-overlapping} outcomes (no outcome counted \ans{more than once}).

\vspace{8pt}
\textbf{VAR-4.A.2 --- Classical Probability Formula} (equally likely outcomes):
\[
  P(E) = \frac{\color{keyred}\textbf{number of outcomes in } \mathbf{E}}{\color{keyred}\textbf{total number of outcomes in } \mathbf{S}}
\]

\vspace{6pt}
\textbf{VAR-4.A.3 --- Probability Range:} For any event $E$:
\ans{$0$} $\le P(E) \le$ \ans{$1$}

\vspace{10pt}
\textbf{Worked Example:} Roll one fair die. $S = \{1, 2, 3, 4, 5, 6\}$
\vspace{4pt}
\renewcommand{\arraystretch}{1.6}
\begin{tabularx}{\linewidth}{@{} l X @{}}
Event & Probability (show work) \\
\hline
$E_1$: roll a 5 & $P(E_1) = $ \ans{$1/6 \approx 0.167$} \\
$E_2$: roll an even number & $P(E_2) = $ \ans{$3/6 = 1/2$} \\
$E_3$: roll a number greater than 4 & $P(E_3) = $ \ans{$2/6 = 1/3$} \\
\end{tabularx}
\end{notesbox}

\vspace{0.1in}

\begin{notesbox}{The Complement Rule}
\small
\textbf{VAR-4.A.4 --- Complement:}\\
The complement of $E$ (written $E^c$ or $E'$) is the event that $E$ \ans{does not occur}.\\[6pt]
\textbf{Complement Rule:} $P(E^c) = $ \ans{$1 - P(E)$}\\[6pt]
\textit{Why it works:} $P(E) + P(E^c) = $ \ans{$1$} because $E$ and $E^c$ together
include \ans{every outcome} in $S$.

\vspace{10pt}
\textbf{Worked Example:} Roll one fair die.
\begin{itemize}[itemsep=4pt]
  \item $P(\text{roll a 5}) = 1/6$, so $P(\text{NOT a 5}) = $ \ans{$1 - 1/6 = 5/6$}
  \item $P(\text{even}) = 1/2$, so $P(\text{odd}) = $ \ans{$1 - 1/2 = 1/2$}
  \item $P(\text{roll} > 4) = 1/3$, so $P(\text{roll} \le 4) = $ \ans{$1 - 1/3 = 2/3$}
\end{itemize}
\end{notesbox}

\vspace{0.1in}

\begin{notesbox}{Interpreting Probability --- Long Run}
\small
\textbf{VAR-4.B.1:} Probabilities of events in \ans{repeatable random} situations can be
interpreted as the relative frequency with which the event will occur \ans{in the long run}.

\vspace{8pt}
\textbf{How to write a long-run interpretation:}\\
``If this random process is repeated many times, the event \ans{(in context)} will occur in
approximately \ans{$P(E) \times 100\%$} of the repetitions.''

\vspace{10pt}
\textbf{Practice:} Write a long-run interpretation for each:
\begin{enumerate}[leftmargin=1.5em, itemsep=10pt]
  \item $P(\text{rolling a 5}) = 1/6$:\\
        \ansline{If a fair die is rolled many times, about 1 in 6 rolls will be a 5.}
  \item A weather forecast says $P(\text{rain tomorrow}) = 0.40$:\\
        \ansline{Over many days with conditions like tomorrow's, about 40\% have rain.}
\end{enumerate}
\end{notesbox}

\vspace{0.1in}

\begin{practicebox}
\small Two dice are rolled. The $6 \times 6$ outcome table is provided.

\vspace{4pt}
\renewcommand{\arraystretch}{1.3}
\begin{tabularx}{\linewidth}{@{} c | *{6}{>{\centering\arraybackslash}X} @{}}
\textbf{Die 1 $\setminus$ Die 2} & \textbf{1} & \textbf{2} & \textbf{3} & \textbf{4} & \textbf{5} & \textbf{6} \\
\hline
\textbf{1} & (1,1) & (1,2) & (1,3) & (1,4) & (1,5) & (1,6) \\
\textbf{2} & (2,1) & (2,2) & (2,3) & (2,4) & (2,5) & (2,6) \\
\textbf{3} & (3,1) & (3,2) & (3,3) & (3,4) & (3,5) & (3,6) \\
\textbf{4} & (4,1) & (4,2) & (4,3) & (4,4) & (4,5) & (4,6) \\
\textbf{5} & (5,1) & (5,2) & (5,3) & (5,4) & (5,5) & (5,6) \\
\textbf{6} & (6,1) & (6,2) & (6,3) & (6,4) & (6,5) & (6,6) \\
\end{tabularx}

\vspace{6pt}
\begin{enumerate}[leftmargin=1.5em, itemsep=8pt]
  \item Total outcomes: \ans{$6 \times 6 = 36$}
  \item $P(\text{sum} = 7)$: circle the outcomes, then calculate. $=$ \ans{$6/36 = 1/6$}
  \item $P(\text{sum} \geq 10) = $ \ans{$6/36 = 1/6$}
  \item $P(\text{NOT sum} \geq 10)$ using complement: $=$ \ans{$1 - 1/6 = 5/6$}
  \item Interpret $P(\text{sum} = 7) = 1/6$ in one sentence:
        \ansline{If two dice are rolled many times, about 1 in 6 rolls sum to 7.}
\end{enumerate}
\end{practicebox}

\end{document}
